\chapter{Software Implementation}

\section{Repository Structure}

The project is implemented in Python with a modular structure.
The main libraries are NumPy, pandas, SciPy, scikit-learn, matplotlib, and PyTorch.

\begin{projectlistingbox}[shellstyle]
src/nn_option_pricing/
    black_scholes.py          # analytical formula and payoff
    config.py                 # experiment configuration dataclasses
    dataset.py                # synthetic data generation and storage
    evaluation.py             # regression metrics
    model.py                  # PyTorch feed-forward neural network
    monte_carlo.py            # efficient Monte Carlo simulation
    noise.py                  # controlled target-noise generation
    noisy_experiment.py       # neural-network noisy-target experiment
    noisy_svr_experiment.py   # noisy-target SVR benchmark
    pipeline.py               # end-to-end experiment
    plots.py                  # visualizations
    svr.py                    # reduced-scale SVR benchmark
    training.py               # training loop, scaling, early stopping
scripts/
    run_experiment.py             # main experiment CLI
    benchmark_runtime.py          # runtime benchmark CLI
    run_svr_benchmark.py          # SVR benchmark CLI
    run_noisy_targets_experiment.py
    run_noisy_svr_benchmark.py
\end{projectlistingbox}
\projectlistingcaption{Main project structure.}

The reusable scientific and machine learning logic is contained in the package \path{src/nn_option_pricing/}.
The files in \texttt{scripts/} are command-line entry points that build configurations and call functions from the package.

\section{Configuration System}

The project uses dataclasses to keep experiment parameters explicit.
The configuration includes dataset ranges, training parameters, Monte Carlo parameters, and output paths.
This avoids hard-coded constants scattered across the codebase and makes experiments easier to reproduce.

Each experiment saves a configuration snapshot in JSON format.
This makes it possible to recover the exact settings used to produce a given result.

\section{Experimental Pipeline}

The main pipeline performs the following steps:

\begin{enumerate}
    \item create output directories;
    \item generate the synthetic dataset;
    \item save the dataset and configuration;
    \item split the data into training, validation, and test sets;
    \item fit input and target scalers;
    \item train the neural network with early stopping;
    \item predict on the held-out test set;
    \item compute neural-network metrics;
    \item run the Monte Carlo benchmark;
    \item save metrics and diagnostic figures.
\end{enumerate}

The final experiment is run through:

\begin{projectlistingbox}[shellstyle]
.venv/bin/python scripts/run_experiment.py \
    --n-samples 100000 \
    --max-epochs 200 \
    --batch-size 1024 \
    --mc-n-paths 50000 \
    --mc-evaluation-samples 512 \
    --feature-set with_moneyness \
    --activation silu \
    --seed 42 \
    --data-dir data/final_improved \
    --output-dir outputs/final_improved
\end{projectlistingbox}
\projectlistingcaption{Command used for the final experiment.}

\section{Testing and Documentation}

The test suite covers:

\begin{itemize}
    \item analytical Black-Scholes pricing;
    \item synthetic dataset generation;
    \item model construction;
    \item Monte Carlo pricing;
    \item target-noise generation;
    \item SVR benchmarks;
    \item end-to-end pipeline smoke tests.
\end{itemize}

The project also includes Sphinx documentation generated from narrative documentation and Python docstrings.
This supports code readability and makes the implementation easier to inspect.

\section{Computational Efficiency}

The Black-Scholes formula and dataset generation are vectorized with NumPy.
The Monte Carlo method uses the exact terminal distribution of the geometric Brownian motion, avoiding unnecessary time discretization.
Monte Carlo simulation is batched both over options and over simulated paths, which keeps memory usage under control.

The neural network training uses PyTorch \texttt{DataLoader}, mini-batch gradient descent, input standardization, and target scaling.
Prediction is performed in batches under \texttt{torch.no\_grad()}.
These choices make the code suitable both for quick development runs and for larger experiments.
